\section{The \texttt{Union\_logger\_1D} McStas Component}
One dimensional Union logger for several possible variables

\subsection*{Identification}
\begin{itemize}
  \item \textbf{Author:} Mads Bertelsen
  \item \textbf{Origin:} University of Copenhagen
  \item \textbf{Date:} 20.08.15
\end{itemize}

\subsection*{Description}
Part of the Union components, a set of components that work together and thus sperates geometry and physics within McStas. The use of this component requires other components to be used.

1) One specifies a number of processes using process components 2) These are gathered into material definitions using Union\_make\_material 3) Geometries are placed using Union\_box/cylinder/sphere, assigned a material 4) A Union\_master component placed after all of the above

Only in step 4 will any simulation happen, and per default all geometries defined before this master, but after the previous will be simulated here.

There is a dedicated manual available for the Union\_components

This logger in particular is not finished. It is supposed to allow several different variables to be histogrammed, but currently only time is supported

A logger will log something for scattering events happening to certain volumes, which are specified in the target\_geometry string. By leaving it blank, all geometries are logged, even the ones not defined at this point in the instrument file. If a list og target\_geometries is selected, one can further narrow the events logged by providing a list of process names in target\_process which need to correspond with names of defined Union\_process components.

To use the logger\_conditional\_extend function, set it to some integer value n and make and extend section to the master component that runs the geometry. In this extend function, logger\_conditional\_extend[n] is 1 if the conditional stack evaluated to true, 0 if not. This way one can check what rays is logged using regular McStas monitors. Only works if a conditional is applied to this logger.

\subsection*{Input parameters}
Parameters in \textbf{boldface} are required; the others are optional.

\begin{longtable}{p{0.22\textwidth}p{0.12\textwidth}p{0.46\textwidth}p{0.14\textwidth}}
\toprule
\textbf{Name} & \textbf{Unit} & \textbf{Description} & \textbf{Default} \\
\midrule
\endhead
target\_geometry & string & Comma seperated list of geometry names that will be logged, leave empty for all volumes (even not defined yet) & "NULL" \\
target\_process & string & Comma seperated names of physical processes, if volumes are selected, one can select Union\_process names & "NULL" \\
\textbf{min\_value} & 1 & Histogram boundery for logged value &  \\
\textbf{max\_value} & 1 & Histogram boundery for logged value &  \\
n1 & 1 & Number of bins in histogram & 90 \\
variable & string & Time for time in seconds, q for magnitude of scattering vector in 1/\AA{}. & "time" \\
filename & string & Filename of produced data file & "NULL" \\
order\_total & 1 & Only log rays that scatter for the n'th time, 0 for all orders & 0 \\
order\_volume & 1 & Only log rays that scatter for the n'th time in the same geometry & 0 \\
order\_volume\_process & 1 & Only log rays that scatter for the n'th time in the same geometry, using the same process & 0 \\
logger\_conditional\_extend\_index & 1 & If a conditional is used with this logger, the result of each conditional calculation can be made available in extend as a array called "logger\_conditional\_extend", and one would then acces logger\_conditional\_extend[n] if logger\_conditional\_extend\_index is set to n & -1 \\
init & string & Name of Union\_init component (typically "init", default) & "init" \\
nowritefile & 1 & If set, logger will skip writing to disk & 0 \\
\bottomrule
\end{longtable}

\subsection*{Links}
\begin{itemize}
  \item Component source code found in file \texttt{Union\_logger\_1D.comp}.
\end{itemize}
\IfFileExists{union/Union_logger_1D_static.tex}{\input{union/Union_logger_1D_static.tex}}{}